% ===================================================================
% SL-05d-ML: Stundenleistung MUSTERLÖSUNG
% Spanisch Klasse 7 - Sequenz 5: Mi tiempo libre
% ===================================================================

\documentclass[11pt, a4paper]{article}

\newif\ifswmodus
\swmodusfalse

\usepackage[utf8]{inputenc}
\usepackage[T1]{fontenc}
\usepackage[ngerman]{babel}
\usepackage{helvet}
\renewcommand{\familydefault}{\sfdefault}

\usepackage[margin=2cm]{geometry}
\usepackage{enumitem}
\usepackage{tabularx}
\usepackage{booktabs}

\usepackage{xcolor}
\usepackage{tcolorbox}
\tcbuselibrary{skins}

\usepackage{fancyhdr}
\usepackage{lastpage}
\usepackage{needspace}

\definecolor{spanisch-color}{HTML}{c41e3a}
\definecolor{grau-color}{HTML}{888888}
\definecolor{hellgrau-color}{HTML}{f5f5f5}
\definecolor{loesung-color}{HTML}{c62828}

\definecolor{sw-dark}{HTML}{000000}
\definecolor{sw-gray}{HTML}{666666}
\definecolor{sw-white}{HTML}{ffffff}

\ifswmodus
    \colorlet{spanisch}{sw-dark}
    \colorlet{grau}{sw-gray}
    \colorlet{hellgrau}{sw-white}
    \colorlet{loesung}{sw-dark}
\else
    \colorlet{spanisch}{spanisch-color}
    \colorlet{grau}{grau-color}
    \colorlet{hellgrau}{hellgrau-color}
    \colorlet{loesung}{loesung-color}
\fi

\newcommand{\fachfarbe}{spanisch}

\newcommand{\thema}{Stundenleistung: Mi tiempo libre}
\newcommand{\fach}{Spanisch}
\newcommand{\klasse}{7}
\newcommand{\maxpunkte}{20}
\newcommand{\zeit}{25 Minuten}
\newcommand{\datum}{\today}

\pagestyle{fancy}
\fancyhf{}
\renewcommand{\headrulewidth}{0.4pt}
\renewcommand{\footrulewidth}{0.4pt}

\lhead{\fach{} -- Klasse \klasse}
\chead{\textcolor{loesung}{\textbf{SL-05d MUSTERLÖSUNG}}}
\rhead{\datum}

\lfoot{\zeit}
\cfoot{}
\rfoot{Seite \thepage\ von \pageref{LastPage}}

\newcommand{\punktebox}[1]{\hfill\fbox{\small \_/#1}}
\newcommand{\luecke}[1]{\rule{#1}{0.4pt}}
\newcommand{\lsg}[1]{\textcolor{loesung}{\textbf{#1}}}

\newtcolorbox{infobox}{
    colback=hellgrau,
    colframe=\fachfarbe,
    boxrule=1pt,
    arc=0pt,
    left=8pt, right=8pt, top=8pt, bottom=8pt
}

\newenvironment{aufgabe}[1]{%
    \needspace{5\baselineskip}%
    \nopagebreak[4]%
    \vspace{10pt}
    \noindent\textbf{\textcolor{\fachfarbe}{Aufgabe #1}}
    \nopagebreak[4]%
    \vspace{5pt}
    \nopagebreak[4]%
    \begin{list}{}{\leftmargin=0pt}
    \item[]
}{%
    \end{list}
}

\newcommand{\es}[1]{\textcolor{\fachfarbe}{\textbf{#1}}}

\begin{document}

\begin{center}
    {\Large\bfseries\textcolor{\fachfarbe}{\thema}}\\[0.3em]
    {\large\textcolor{loesung}{\textbf{--- MUSTERLÖSUNG ---}}}
\end{center}

\vspace{5pt}

\noindent
\begin{tabularx}{\textwidth}{|X|X|X|}
\hline
\textbf{Name:} \lsg{Musterschüler/in} &
\textbf{Klasse:} \klasse &
\textbf{Datum:} \datum \\
\hline
\end{tabularx}

\vspace{10pt}

\begin{infobox}
\textbf{Zeit:} \zeit{} | \textbf{Punkte:} \maxpunkte{} | \textbf{Hilfsmittel:} keine
\end{infobox}

\vspace{15pt}

% ===================================================================
% AUFGABE 1
% ===================================================================

\begin{aufgabe}{1 -- Hobbys zuordnen}
Welches Verb passt? Schreibe \es{jugar}, \es{tocar} oder \es{hacer}. \punktebox{4}
\end{aufgabe}

\noindent
a) \lsg{jugar} al fútbol \hfill
b) \lsg{tocar} la guitarra \\[0.8em]
c) \lsg{hacer} deporte \hfill
d) \lsg{jugar} al baloncesto \\

\lsg{Bewertung: je 1 Punkt}

\vspace{15pt}

% ===================================================================
% AUFGABE 2
% ===================================================================

\begin{aufgabe}{2 -- Verben konjugieren}
Ergänze die richtige Form des Verbs. \punktebox{6}
\end{aufgabe}

\noindent
a) Yo \lsg{como} (comer) una pizza. \hfill (ich esse)

\vspace{0.8em}

\noindent
b) Tú \lsg{bebes} (beber) agua. \hfill (du trinkst)

\vspace{0.8em}

\noindent
c) Ella \lsg{vive} (vivir) en Madrid. \hfill (sie lebt)

\vspace{0.8em}

\noindent
d) Nosotros \lsg{leemos} (leer) un libro. \hfill (wir lesen)

\vspace{0.8em}

\noindent
e) Él \lsg{escribe} (escribir) una carta. \hfill (er schreibt)

\vspace{0.8em}

\noindent
f) Ellos \lsg{comen} (comer) en casa. \hfill (sie essen)

\lsg{Bewertung: je 1 Punkt}

\vspace{15pt}

% ===================================================================
% AUFGABE 3
% ===================================================================

\begin{aufgabe}{3 -- Sätze mit gustar}
Vervollständige die Sätze mit \es{me gusta}, \es{te gusta} oder \es{le gusta}. \punktebox{4}
\end{aufgabe}

\noindent
a) A mí \lsg{me gusta} escuchar música. \hfill (Ich höre gern Musik.)

\vspace{0.8em}

\noindent
b) A ti \lsg{te gusta} jugar a videojuegos. \hfill (Du spielst gern Videospiele.)

\vspace{0.8em}

\noindent
c) A él \lsg{le gusta} leer libros. \hfill (Er liest gern Bücher.)

\vspace{0.8em}

\noindent
d) A mí \lsg{me gusta} hacer deporte. \hfill (Ich mache gern Sport.)

\lsg{Bewertung: je 1 Punkt}

\vspace{15pt}

% ===================================================================
% AUFGABE 4
% ===================================================================

\begin{aufgabe}{4 -- Meine Hobbys}
Schreibe 3 Sätze über deine Hobbys auf Spanisch. \punktebox{6}
\end{aufgabe}

\noindent
\textbf{Verwende:} Me gusta... / Me llamo... / Vivo en...

\vspace{0.8em}

\noindent
1. \lsg{Me llamo María. (oder anderer Name)}

\vspace{1em}

\noindent
2. \lsg{Me gusta jugar al fútbol. (oder anderes Hobby)}

\vspace{1em}

\noindent
3. \lsg{Vivo en Rostock. (oder anderer Ort)}

\vspace{1em}

\lsg{Bewertung:}
\begin{itemize}
    \item 2 Punkte pro Satz (1P Grammatik + 1P Inhalt/Wortschatz)
    \item Akzeptierte Strukturen: Me llamo..., Vivo en..., Me gusta + Infinitiv
    \item Kleine Rechtschreibfehler: halber Punkt Abzug
\end{itemize}

% ===================================================================
% PUNKTEÜBERSICHT
% ===================================================================

\vspace{25pt}
\hrule
\vspace{10pt}

\noindent
\begin{tabularx}{\textwidth}{|X|c|c|c|c||c|}
\hline
\textbf{Aufgabe} & 1 & 2 & 3 & 4 & \textbf{Gesamt} \\
\hline
\textbf{max. Punkte} & 4 & 6 & 4 & 6 & \textbf{20} \\
\hline
\textbf{erreicht} & \lsg{4} & \lsg{6} & \lsg{4} & \lsg{6} & \lsg{20} \\
\hline
\end{tabularx}

\vspace{10pt}

\begin{flushright}
\textbf{Note:} \lsg{1}
\end{flushright}

\vspace{1em}

\textbf{Notenschlüssel:}

\begin{tabular}{|c|c|c|c|c|c|c|}
\hline
\textbf{Note} & 1 & 2 & 3 & 4 & 5 & 6 \\
\hline
\textbf{ab Punkte} & 19 & 16 & 12 & 8 & 4 & <4 \\
\hline
\end{tabular}

\end{document}
